% Vietnamese translation for OI Public Library
% Source-PDF: 动态规划 DYNAMIC PROGRAMMING/楼天城牛关于队列优化多重背包.pdf
\documentclass[11pt]{article}
\usepackage{oipl-vn}

\newcommand{\Oh}{\mathcal{O}}
\newcommand{\Max}{\operatorname{max}}
\newcommand{\Min}{\operatorname{min}}

\title{Một số ví dụ quy hoạch động phân phối tài nguyên}
\author{Lou Tiancheng}
\date{}

\begin{document}
\maketitle

\origtitle{Resource-allocation dynamic programming examples}
\sourcepath{Dynamic Programming / Queue-optimized multiple knapsack - Lou Tiancheng}

\begin{abstract}
Tài liệu nguồn là một bộ slide ngắn về các bài quy hoạch động dạng phân phối
tài nguyên và cái túi: phân phối máy, độ tin cậy hệ thống, sản xuất suất ăn
nhanh, mua hàng có ưu đãi và bài ngân sách của Jinming. Bản dịch chuyển nội
dung slide thành ghi chú liền mạch, nhấn mạnh trạng thái, chuyển trạng thái,
khởi tạo và các cách giảm phạm vi lặp.
\end{abstract}

\tableofcontents

\section{Phân phối máy}

\subsection{Đề bài}

Công ty mẹ có \(M\) máy sản xuất hiệu suất cao và muốn phân phối cho \(N\)
công ty con. Nếu công ty con \(i\) nhận \(j\) máy, lợi nhuận thu được là
\(\operatorname{value}[i,j]\). Mỗi công ty có thể nhận số máy tùy ý, nhưng
tổng số máy được cấp không vượt quá \(M\). Ràng buộc trong slide là
\[
  M\le 15,\qquad N\le 10.
\]
Dữ liệu gồm \(M,N\), sau đó là ma trận lợi nhuận kích thước \(N\times M\).

\subsection{Phân tích}

Dùng số máy làm tài nguyên. Đặt
\[
  F[i,j]
\]
là lợi nhuận lớn nhất khi phân phối \(j\) máy cho \(i\) công ty đầu. Nếu công
ty \(i\) nhận \(j-k\) máy, \(k\) máy còn lại đã được phân phối cho \(i-1\)
công ty trước đó. Công thức:
\[
  F[i,j]=\max_{0\le k\le j}
  \{F[i-1,k]+\operatorname{value}[i,j-k]\}.
\]
Điều kiện đầu là \(F[0,0]=0\). Độ phức tạp là \(\Oh(NM^2)\).

\section{Độ tin cậy hệ thống}

\subsection{Đề bài}

Một hệ thống gồm nhiều bộ phận mắc nối tiếp. Chỉ cần một bộ phận hỏng thì cả
hệ thống không hoạt động. Để tăng độ tin cậy, mỗi bộ phận có thể lắp thêm các
linh kiện dự phòng; càng nhiều dự phòng thì xác suất hoạt động càng cao nhưng
chi phí cũng tăng.

Với bộ phận \(i\), biết đơn giá dự phòng \(\operatorname{cost}[i]\) và xác
suất hoạt động \(P[i,k]\) khi dùng \(k\) linh kiện dự phòng. Tổng chi phí bị
giới hạn bởi \(C\). Hãy tối đa hóa xác suất hệ thống hoạt động.

\subsection{Phân tích}

Vì các bộ phận mắc nối tiếp, xác suất hệ thống hoạt động là tích xác suất của
các bộ phận. Đặt \(F[i,m]\) là độ tin cậy lớn nhất khi dùng ngân sách \(m\)
cho \(i\) bộ phận đầu. Khi bộ phận \(i\) dùng \(k\) linh kiện dự phòng:
\[
  F[i,m]=
  \max_{0\le k\le \lfloor m/\operatorname{cost}[i]\rfloor}
  F[i-1,m-k\operatorname{cost}[i]]\cdot P[i,k].
\]
Đáp án là \(F[n,C]\). Trong cài đặt thực tế, cần khởi tạo trạng thái xác suất
theo nghĩa ``không có bộ phận nào'' là xác suất \(1\) tại chi phí \(0\); các
trạng thái không đạt được đặt thành giá trị không hợp lệ.

\section{Bài toán suất ăn nhanh}

\subsection{Đề bài}

Một suất ăn gồm \(A\) hamburger, \(B\) phần khoai và \(C\) đồ uống. Thời gian
sản xuất một đơn vị của ba loại lần lượt là \(p_1,p_2,p_3\). Có \(N\) dây
chuyền sản xuất, dây chuyền \(i\) có thời gian làm việc \(T[i]\). Mỗi dây
chuyền có thể sản xuất cả ba loại sản phẩm. Hỏi trong một ngày có thể sản xuất
nhiều nhất bao nhiêu suất ăn. Giả sử sản lượng mỗi loại không vượt quá \(100\).

\subsection{Quy hoạch động}

Vì các dây chuyền độc lập, dùng dây chuyền làm giai đoạn. Đặt
\[
  p[i,j,k]
\]
là số đồ uống lớn nhất có thể sản xuất sau khi dùng \(i\) dây chuyền đầu,
trong đó đã sản xuất \(j\) hamburger và \(k\) phần khoai. Đặt
\[
  r[i,x,y]
\]
là số đồ uống nhiều nhất mà riêng dây chuyền \(i\) có thể sản xuất nếu trên
dây chuyền này đã sản xuất \(x\) hamburger và \(y\) phần khoai. Khi đó
\[
  r[i,x,y]=
  \left\lfloor
  \frac{T[i]-xp_1-yp_2}{p_3}
  \right\rfloor
\]
nếu \(xp_1+yp_2\le T[i]\), và trạng thái không hợp lệ nếu ngược lại.

Chuyển trạng thái:
\[
  p[i,j,k]=
  \max_{\substack{0\le j_1\le j\\0\le k_1\le k}}
  \{p[i-1,j_1,k_1]+r[i,j-j_1,k-k_1]\}.
\]
Cách trực tiếp có độ phức tạp \(\Oh(N\cdot 100^4)\).

\subsection{Cắt tỉa và tối ưu}

Ta có thể tính trước một cận trên hiển nhiên cho số suất ăn:
\[
  \min\left(
    \left\lfloor\frac{100}{A}\right\rfloor,
    \left\lfloor\frac{100}{B}\right\rfloor,
    \left\lfloor\frac{100}{C}\right\rfloor
  \right).
\]
Sau đó dùng một phương án tham lam nhanh để lấy một nghiệm khả thi. Nếu nghiệm
khả thi đã đạt cận trên thì có thể kết thúc ngay. Trong quá trình DP, sau mỗi
giai đoạn cũng có thể so với cận trên để dừng sớm.

Hai tối ưu kỹ thuật khác trong slide:
\begin{itemize}
  \item Vì giai đoạn \(i\) chỉ phụ thuộc giai đoạn \(i-1\), dùng hai mảng
  \(100\times100\) và hoán đổi con trỏ thay vì sao chép toàn bộ mảng.
  \item Thu hẹp biên vòng lặp bằng các cận trên \(q_1,q_2\) cho số hamburger
  và khoai, đồng thời giới hạn \(j_1,k_1\) theo thời gian còn lại của dây
  chuyền hiện tại và tổng thời gian của các dây chuyền trước.
\end{itemize}

\section{Mua hàng có ưu đãi}

\subsection{Đề bài}

Một cửa hàng bán tối đa \(5\) loại hàng trong giỏ của khách, mỗi loại mua
không quá \(5\) món. Ngoài giá đơn lẻ, cửa hàng còn có tối đa \(99\) gói ưu
đãi; mỗi gói gồm một số lượng cụ thể của một hoặc nhiều loại hàng và có giá
thấp hơn mua lẻ. Cần tính số tiền nhỏ nhất để mua đúng số hàng trong giỏ,
không được mua thừa dù mua thừa có thể rẻ hơn.

Ví dụ: một bông hoa giá \(2\), một bình hoa giá \(5\). Ưu đãi ba bông hoa giá
\(5\), hoặc hai bình hoa cộng một bông hoa giá \(10\). Nếu cần mua ba bông hoa
và hai bình hoa, giá tối ưu là \(14\): dùng gói một bông hoa cộng hai bình
hoa giá \(10\), rồi mua lẻ hai bông hoa giá \(4\).

\subsection{Phân tích}

Vì có nhiều nhất \(5\) loại hàng và mỗi loại cần nhiều nhất \(5\), có thể dùng
trạng thái năm chiều:
\[
  F[a,b,c,d,e]
\]
là chi phí nhỏ nhất để mua các số lượng tương ứng của năm loại hàng. Trạng
thái khởi tạo chính là mua lẻ:
\[
  F[a,b,c,d,e]=
  a\operatorname{cost}_1+b\operatorname{cost}_2+
  c\operatorname{cost}_3+d\operatorname{cost}_4+e\operatorname{cost}_5.
\]

Với mỗi gói ưu đãi \(s\), ký hiệu số lượng của năm loại hàng trong gói là
\((s_1,s_2,s_3,s_4,s_5)\), giá gói là \(\operatorname{saleCost}[s]\). Nếu
trạng thái hiện tại đủ hàng để dùng ngược lại gói này, ta cập nhật:
\[
\begin{aligned}
F[a,b,c,d,e]=\min\{F[a,b,c,d,e],\;&
F[a-s_1,b-s_2,c-s_3,d-s_4,e-s_5]\\
&+\operatorname{saleCost}[s]\}.
\end{aligned}
\]

Bài toán thỏa tính không hậu nghiệm: sau khi dùng một gói ưu đãi, phần còn lại
là cùng một kiểu bài toán với số lượng nhỏ hơn. Vì gói ưu đãi được dùng không
giới hạn số lần, ta chỉ cần lặp qua toàn bộ không gian trạng thái nhỏ.

\section{Bài ngân sách của Jinming}

\subsection{Đề bài}

Jinming muốn mua đồ cho phòng của mình trong ngân sách \(N\). Mỗi đồ vật có
giá \(v[j]\), độ quan trọng \(w[j]\) từ \(1\) đến \(5\), và có thể là vật
chính hoặc phụ kiện của một vật chính. Nếu mua phụ kiện thì bắt buộc phải mua
vật chính tương ứng. Mỗi vật chính có \(0,1\) hoặc \(2\) phụ kiện; phụ kiện
không có phụ kiện con. Cần tối đa hóa
\[
  \sum v[j]\cdot w[j]
\]
trong ngân sách.

\subsection{Nếu chỉ có vật chính}

Khi không có phụ kiện, đây là cái túi 01 chuẩn. Đặt \(F[i,j]\) là giá trị lớn
nhất khi xét \(i\) vật phẩm đầu và dùng không quá \(j\) tiền:
\[
  F[i,j]=\max\{F[i-1,j],F[i-1,j-C_i]+V_i\}.
\]
Độ phức tạp là \(\Oh(nm)\), với \(n\) là ngân sách và \(m\) là số vật phẩm.

\subsection{Quay lại bài gốc}

Với một vật chính \(i\), ta liệt kê các phương án phụ kiện hợp lệ. Nếu vật
chính có hai phụ kiện thì tối đa có bốn phương án: chỉ vật chính, vật chính
kèm phụ kiện thứ nhất, vật chính kèm phụ kiện thứ hai, hoặc vật chính kèm cả
hai. Mỗi phương án trở thành một ``vật phẩm tổng hợp'' có chi phí và giá trị
bằng tổng của các món trong phương án.

Gọi các phương án của vật chính \(i\) có chi phí \(C_{i,k}\) và giá trị
\(V_{i,k}\). Công thức:
\[
  F[i,j]=
  \max\left\{
    F[i-1,j],\;
    F[i-1,j-C_{i,k}]+V_{i,k}
  \right\}.
\]
Trong đó \(k\) chạy qua các phương án phụ kiện hợp lệ. Có thể thêm các phụ
kiện ảo có chi phí và giá trị \(0\) để mỗi vật chính luôn có đúng hai phụ
kiện, rồi liệt kê chọn hoặc không chọn từng phụ kiện. Độ phức tạp vẫn là
\(\Oh(nm)\) vì số phương án của mỗi vật chính là hằng số.

\section*{Kết luận}

Các ví dụ trong tài liệu đều có cùng khuôn mẫu: chọn một tài nguyên làm chiều
trạng thái, chia quá trình thành giai đoạn độc lập, rồi liệt kê phần tài
nguyên cấp cho giai đoạn hiện tại. Khi số trạng thái nhỏ, cách liệt kê trực
tiếp đủ tốt; khi không gian lớn, cần dùng mảng cuộn, cận trên, cắt tỉa vòng
lặp hoặc đóng gói nhiều lựa chọn thành một nhóm vật phẩm.

\end{document}
